
\SGARefTarget{sga3:xxiv:2.9}
\medskip\noindent\textbf{Remark 2.9.} ---
It is not true in general that \(B\) and \(B'\) are isomorphic;
they are nevertheless inner twisted forms of one another
(cf.\ \S5).

One may develop ``\(\operatorname{Isomint}\)'' variants
of the preceding results. We mention one:

\SGARefTarget{sga3:xxiv:2.10}
\medskip\noindent\textbf{Proposition 2.10.} ---
{\itshape Let \(S\) be a scheme, let \(G\) and \(G'\) be two reductive
\(S\)-groups having the same type at every point, let \(B\supset T\)
(respectively, \(B'\supset T'\)) be a Killing couple of \(G\)
(respectively, \(G'\)), and let
\(u\in\underline{\operatorname{Isomext}}(G,G')(S)\). Consider}
\[
\begin{aligned}
  &\underline{\operatorname{Isomint}}_u
     (G,B,T;G',B',T')\\
  &\qquad =
  \underline{\operatorname{Isom}}_{S\text{-}\mathrm{gr}}
     (G,B,T;G',B',T')
  \cap
  \underline{\operatorname{Isomint}}_u(G,G').
\end{aligned}
\]
{\itshape It is a smooth affine \(S\)-scheme which is a principal
homogeneous bundle under \(T_{\mathrm{ad}}\). In particular,}
\[
  \underline{\operatorname{Isomint}}_u
    (G,B,T;G',B',T')(S)\ne\varnothing
\]
{\itshape if and only if the corresponding element of
\(H^1(S,T_{\mathrm{ad}})\) is zero.}

To conclude this section, we prove two results that will be useful
later.

\SGARefTarget{sga3:xxiv:2.11}
\medskip\noindent\textbf{Proposition 2.11.} ---
{\itshape Let \(S\) be a scheme, let \(G\) be a reductive \(S\)-group,
and let \(T\) be a maximal torus of \(G\). The evident morphism}
\[
  T_{\mathrm{ad}}
  \xrightarrow{\ \sim\ }
  \underline{\Aut}_{S\text{-}\mathrm{gr}}(G,\operatorname{id}_T)
\]
{\itshape is an isomorphism.}

This follows from the preceding statements. A direct proof was also
given in the course of the proof of \SGARefLink{sga3:xxiv:1.5.2}{1.5.2}.

\SGARefTarget{sga3:xxiv:2.12}
\medskip\noindent\textbf{Corollary 2.12.} ---
{\itshape Under the preceding hypotheses, there is an equivalence
between the category of pairs \((G',f)\), where \(G'\) is a form of
\(G\) and \(f\) is an isomorphism from \(T\) onto a maximal torus of
\(G'\), and the category of principal homogeneous bundles under
\(T_{\mathrm{ad}}\).}

\SGARefTarget{sga3:xxiv:2.13}
\medskip\noindent\textbf{Corollary 2.13.} ---
{\itshape If \(H^1(S,T_{\mathrm{ad}})=0\), and if \(G'\) is a form of
\(G\) having a maximal torus isomorphic to \(T\), then \(G'\) is
isomorphic to \(G\).}

\SGARefTarget{sga3:xxiv:2.14}
\medskip\noindent\textbf{Corollary 2.14.} ---
{\itshape Let \(S\) be a scheme such that \(\operatorname{Pic}(S)=0\),
and let \(G\) be a reductive \(S\)-group of constant type. Then \(G\)
is splittable if and only if \(G\) has a split maximal torus.}

Let \(G\) be a reductive group and let
\(\operatorname{rad}(G)\) be its radical (\SGARefLink{sga3:xxii:rem:4.3.9}{Exp.~XXII, 4.3.9}).
Since \(\operatorname{rad}(G)\) is central and characteristic in
\(G\), there is a canonical morphism
\[
  q:\underline{\Autext}(G)
  \longrightarrow
  \underline{\Aut}_{S\text{-}\mathrm{gr}}
    \bigl(\operatorname{rad}(G)\bigr).
\]


\SGARefTarget{sga3:xxiv:2.15}
\medskip\noindent\textbf{Proposition 2.15.} ---
{\itshape Let \(G\) be a reductive \(S\)-group. The following sequence
is exact:}
\[
  1\longrightarrow
  \operatorname{ad}(G)
  \longrightarrow
  \underline{\Aut}_{S\text{-}\mathrm{gr}}
    \bigl(G,\operatorname{id}_{\operatorname{rad}(G)}\bigr)
  \xrightarrow{p}
  \underline{\Autext}(G)
  \xrightarrow{q}
  \underline{\Aut}_{S\text{-}\mathrm{gr}}
    \bigl(\operatorname{rad}(G)\bigr),
\]
{\itshape and \(\ker(q)=\operatorname{im}(p)\) is an open and closed
subscheme of \(\underline{\Autext}(G)\), finite over \(S\).}

One may assume that \(G\) is split. The first assertion is immediate,
and the second follows from \SGARefLink{sga3:XXI:6.7.5}{Exp.~XXI, 6.7.5} and 6.7.7.

Writing
\[
  H=
  \underline{\Aut}_{S\text{-}\mathrm{gr}}
    \bigl(G,\operatorname{id}_{\operatorname{rad}(G)}\bigr),
\]
one obtains:

\SGARefTarget{sga3:xxiv:2.16}
\medskip\noindent\textbf{Corollary 2.16.} ---
{\itshape There is an equivalence between the category of pairs
\((G',f)\), where \(G'\) is a form of \(G\) and \(f\) is an
isomorphism from \(\operatorname{rad}(G)\) onto the radical of \(G'\),
and the category of principal bundles under a certain \(S\)-group
scheme \(H\), where \(H\) is such that there is an exact sequence}
\[
  1\longrightarrow \operatorname{ad}(G)
  \longrightarrow H
  \longrightarrow F
  \longrightarrow 1,
\]
{\itshape in which the \(S\)-group \(F\) is étale and finite over
\(S\).}

\SGARefTarget{sga3:xxiv:section:3}
\section*{3. Dynkin Scheme of a Reductive Group. Quasi-Split Groups}

\SGARefTarget{sga3:xxiv:3.1}
\medskip\noindent\textbf{3.1.} ---
Recall (\SGARefLink{sga3:XXI:7.4.1}{Exp.~XXI, 7.4.1}) that a Dynkin diagram is a finite set endowed
with the structure defined by a set of pairs of distinct elements
(bonds) and a map to \(\{1,2,3\}\) (lengths). To every pinned reduced
root datum \(\mathscr R\) there is associated a Dynkin diagram
\(\Delta(\mathscr R)\), whose underlying set is the set of simple
roots.

\SGARefTarget{sga3:xxiv:3.2}
\medskip\noindent\textbf{3.2.} ---
Let \(S\) be a scheme. A Dynkin \(S\)-scheme is a finite twisted
constant \(S\)-scheme \(\Delta\), endowed with the structure defined
by a subscheme \(L\) of \(\Delta\times_S\Delta\) having empty
intersection with the diagonal, and by a morphism
\[
  \Delta\longrightarrow\{1,2,3\}_S.
\]
If \(S'\) is a connected \(S\)-scheme,
\(\Delta(S')\) is naturally endowed with the structure of a Dynkin
diagram.

The following notions are defined immediately: an isomorphism of two
Dynkin schemes, the base extension of a Dynkin scheme, and the
constant Dynkin scheme associated with a Dynkin diagram.

Every étale descent datum on a Dynkin scheme is effective.

\SGARefTarget{sga3:xxiv:3.3}
\medskip\noindent\textbf{3.3.} ---
We shall associate a Dynkin \(S\)-scheme with every reductive
\(S\)-group \(G\). First suppose that \(G\) is splittable over \(S\).
For every pinning \(E\) of \(G\), let \(\Delta(E)\) be the constant
Dynkin scheme associated with the pinned root datum defined by \(E\).
If \(E\) and \(E'\) are two pinnings of \(G\), then by \SGARefLink{sga3:xxiv:1.5}{1.5} there is
a unique inner automorphism of \(G\) over \(S\) carrying \(E\) to
\(E'\). This automorphism of \(G\) defines an isomorphism
\[
  a_{EE'}:\Delta(E)\xrightarrow{\ \sim\ }\Delta(E').
\]
The \(a_{EE'}\) plainly form a transitive system, so the
\(\Delta(E)\) may be identified (that is, one takes their inductive
limit); the result is a constant Dynkin scheme denoted by
\(\operatorname{Dyn}(G)\). If now \(G\) is an arbitrary reductive
\(S\)-group, there is an étale covering family \(\{S_i\to S\}\)
such that


\(G_{S_i}\) is splittable. Arguing as above, one therefore obtains a
canonical descent datum on the \(\operatorname{Dyn}(G_{S_i})\), and
by descent this constructs a Dynkin \(S\)-scheme
\(\operatorname{Dyn}(G)\).

\SGARefTarget{sga3:xxiv:3.4}
\medskip\noindent\textbf{3.4.} ---
This construction has the following properties (which, moreover,
essentially characterize it):
\begin{enumerate}
\item[\SGAInlineRefTarget{sga3:xxiv:3.4:item:i}(i)] With every reductive \(S\)-group \(G\) is associated a
Dynkin scheme \(\operatorname{Dyn}(G)\); with every isomorphism
\[
  u:G\xrightarrow{\ \sim\ }G'
\]
is functorially associated an isomorphism
\[
  \operatorname{Dyn}(u):
  \operatorname{Dyn}(G)\xrightarrow{\ \sim\ }\operatorname{Dyn}(G').
\]

\item[\SGAInlineRefTarget{sga3:xxiv:3.4:item:ii}(ii)] If \(S'\) is an \(S\)-scheme and \(G\) is a reductive
\(S\)-group, then
\[
  \operatorname{Dyn}(G\times_S S')
  \simeq
  \operatorname{Dyn}(G)\times_S S'.
\]

\item[(iii)] If \(E\) is a pinning of \(G\) defining the pinned root
datum with Dynkin diagram \(\Delta\), then
\[
  \operatorname{Dyn}(G)\simeq\Delta_S.
\]

\item[\SGAInlineRefTarget{sga3:xxiv:3.4:item:iv}(iv)] If \(u\) is an inner automorphism of \(G\), then
\(\operatorname{Dyn}(u)\) is the identity automorphism of
\(\operatorname{Dyn}(G)\).
\end{enumerate}

\SGARefTarget{sga3:xxiv:3.5}
\medskip\noindent\textbf{3.5.} ---
Let \(S\) be a scheme and \(G\) a reductive \(S\)-group. It is clear
that the functor
\[
  \underline{\Aut}_{\mathrm{Dyn}}
    \bigl(\operatorname{Dyn}(G)\bigr)
\]
of automorphisms of \(\operatorname{Dyn}(G)\) as a Dynkin scheme is
represented by a finite twisted constant \(S\)-scheme. By \SGARefLink{sga3:xxiv:3.4:item:i}{3.4(i)}
and \SGARefLink{sga3:xxiv:3.4:item:ii}{(ii)}, there is a canonical morphism
\[
  \underline{\Aut}_{S\text{-}\mathrm{gr}}(G)
  \longrightarrow
  \underline{\Aut}_{\mathrm{Dyn}}
    \bigl(\operatorname{Dyn}(G)\bigr),
\]
which, by \SGARefLink{sga3:xxiv:3.4:item:iv}{(iv)}, factors through a morphism
\[
  \underline{\Autext}(G)
  \longrightarrow
  \underline{\Aut}_{\mathrm{Dyn}}
    \bigl(\operatorname{Dyn}(G)\bigr).
\]
More generally, if \(G\) and \(G'\) are two reductive \(S\)-groups,
there is a canonical morphism
\[
  \underline{\operatorname{Isomext}}(G,G')
  \longrightarrow
  \underline{\operatorname{Isom}}_{\mathrm{Dyn}}
    \bigl(\operatorname{Dyn}(G),\operatorname{Dyn}(G')\bigr).
\]
In particular, if \(G'\) is an inner twisted form of \(G\) (\SGARefLink{sga3:xxiv:1.11}{1.11}),
the Dynkin schemes \(\operatorname{Dyn}(G)\) and
\(\operatorname{Dyn}(G')\) are isomorphic.

\SGARefTarget{sga3:xxiv:3.6}
\medskip\noindent\textbf{3.6.} ---
If \(G\) is semisimple (respectively, adjoint or simply connected),
the morphism
\[
  \underline{\Autext}(G)
  \longrightarrow
  \underline{\Aut}_{\mathrm{Dyn}}
    \bigl(\operatorname{Dyn}(G)\bigr)
\]
is a monomorphism (respectively, an isomorphism). Indeed, one may
assume that \(G\) is pinned, and the assertion reduces to the
corresponding result for pinned reduced root data
(cf.\ \SGARefLink{sga3:XXI:7.4.5}{Exp.~XXI, 7.4.5}).

There is an analogous result for the \(\operatorname{Isom}\) functors.
It follows in particular that two semisimple adjoint (respectively,
simply connected) \(S\)-groups are inner twisted forms of one another
if and only if their Dynkin schemes are isomorphic.


\SGARefTarget{sga3:xxiv:3.7}
\medskip\noindent\textbf{3.7.} ---
One can give a different construction of the Dynkin scheme associated
with a reductive group. Let \(\mathscr R\) be a pinned reduced root
datum and \(G\) a reductive \(S\)-group of type \(\mathscr R\), and
write \(\Delta(\mathscr R)\) for the Dynkin diagram defined by
\(\mathscr R\). There is, by \SGARefLink{sga3:xxiv:3.5}{3.5}, a canonical morphism
\[
  A_S(\mathscr R)
  =
  \underline{\Aut}_{S\text{-}\mathrm{gr}}
    \bigl(\operatorname{\acute{E}p}_S(\mathscr R)\bigr)
  \longrightarrow
  \underline{\Aut}_{\mathrm{Dyn}}
    \bigl(\Delta(\mathscr R)_S\bigr).
\]
The reductive \(S\)-group \(G\) corresponds (\SGARefLink{sga3:xxiv:1.17}{1.17}) to a bundle
\[
  \underline{\operatorname{Isom}}_{S\text{-}\mathrm{gr}}
    \bigl(\operatorname{\acute{E}p}_S(\mathscr R),G\bigr),
\]
principal homogeneous under \(A_S(\mathscr R)\). The associated
bundle under
\(\underline{\Aut}_{\mathrm{Dyn}}(\Delta(\mathscr R)_S)\)
corresponds to a form over \(S\) of \(\Delta(\mathscr R)_S\): it is
\(\operatorname{Dyn}(G)\). In other words, this associated bundle is
precisely
\[
  \underline{\operatorname{Isom}}_{\mathrm{Dyn}}
    \bigl(\Delta(\mathscr R)_S,\operatorname{Dyn}(G)\bigr).
\]
In this form the proof is immediate.

\SGARefTarget{sga3:xxiv:3.8}
\medskip\noindent\textbf{3.8. Dynkin Scheme and Killing Couples.} ---
Let \(S\) be a scheme, \(G\) a reductive \(S\)-group, and
\(B\supset T\) a Killing couple of \(G\). There is a canonical
morphism
\[
  i:\operatorname{Dyn}(G)
  \longrightarrow
  \underline{\operatorname{Hom}}_{S\text{-}\mathrm{gr}}
    (T,\mathbb G_{m,S})
\]
which identifies \(\operatorname{Dyn}(G)\) with the ``scheme of
simple roots of \(B\) relative to \(T\).'' This morphism is defined
immediately by descent from the pinned case. Observe also that the
data of \(T\) and \(i\) allow one to reconstruct \(B\) (the
``one-to-one correspondence between systems of simple roots and
systems of positive roots'').

It follows from the preceding description of
\(D=\operatorname{Dyn}(G)\) that there is a canonical root of \(B_D\)
relative to \(T_D\): this root \(\alpha_D\) is the image, under
\(i(D)\), of the identity morphism of \(D\). One obtains from it a
canonical invertible \(\mathcal O_D\)-module \(\mathfrak g_D\):
\[
  \mathfrak g_D
  =
  \bigl(\mathfrak g\otimes_{\mathcal O_S}\mathcal O_D\bigr)^{\alpha_D}.
\]
In the pinned case,
\[
  D=\coprod_{\alpha\in\Delta}S_\alpha,
\]
where every \(S_\alpha\) is a copy of \(S\), and
\(\mathfrak g_D\) is then the \(\mathcal O_D\)-module inducing
\(\mathfrak g_\alpha\) on \(S_\alpha\), for every
\(\alpha\in\Delta\).

\SGARefTarget{sga3:xxiv:3.9}
\medskip\noindent\textbf{3.9. Quasi-Pinnings. Quasi-Pinned Groups.} ---
If \(G\) is a reductive \(S\)-group, a \emph{quasi-pinning of \(G\)}
means the data
\begin{enumerate}
\item[(i)] of a Killing couple \(B\supset T\) of \(G\);
\item[(ii)] of a section
\[
  X\in\Gamma\bigl(\operatorname{Dyn}(G),\mathfrak g_D\bigr)^\times.
\]
\end{enumerate}
A reductive \(S\)-group is called \emph{quasi-splittable} if it has a
quasi-pinning. A \emph{quasi-pinned group} is a reductive group
equipped with a quasi-pinning.

Let \(B\supset T\) be a Killing couple of the reductive \(S\)-group
\(G\). Then \(G\) is quasi-pinnable relative to this Killing couple
if and only if \(\mathfrak g_D\) has a section nonzero at every point,
that is, if the element of \(\operatorname{Pic}(\operatorname{Dyn}(G))\)
defined by \(\mathfrak g_D\) is zero. Suppose in particular that
\(S\) is semilocal. Then \(\operatorname{Dyn}(G)\) is also semilocal,
and hence \(\operatorname{Pic}(\operatorname{Dyn}(G))=0\). Thus:

\SGARefTarget{sga3:xxiv:3.9.1}
\medskip\noindent\textbf{Proposition 3.9.1.} ---
{\itshape Let \(S\) be a semilocal scheme and \(G\) a reductive
\(S\)-group. Then \(G\) is quasi-splittable if and only if it has a
Borel subgroup.}

